\chapter{Conclusion}
\label{ch:conclusion}

\section{What was done}

A complete system was built for driving two worn robot arms and keeping the
person carrying them safe. It has three input routes: a sensed mannequin arm,
a pair of hand controllers used as motion trackers, and ordinary English typed
or spoken. It has a safety layer built on measured distance to the wearer
rather than on collision flags, tested by breaking things rather than by
reading code. It has a perception system that measures the working area
instead of being told about it, and a planner that knows the table is there.
All of it is reachable from one window, and all of it is checked by 1\,322
automated tests and by a program that presses every button in that window.

\section{Against the aims}

\begin{table}[htbp]
  \centering
  \small
  \caption[Aims and outcomes]{The six aims from \cref{ch:intro}, and what
  became of each.}
  \label{tab:aims}
  \begin{tabular}{@{}p{0.42\linewidth}p{0.50\linewidth}@{}}
    \toprule
    \textbf{Aim} & \textbf{Outcome} \\
    \midrule
    Build and characterise the master arm & Done. Six of fourteen channels
      meet the health criterion, and the failing ones are quantified rather
      than described \\
    A mapping that degrades predictably & Done, and measured: one step down
      the ladder costs \SI{63.8}{\milli\metre} of accuracy, and the bottom
      rung is a refusal \\
    A safety layer suited to a worn machine & Built and validated by
      injection, 14 of 14. One defect remains open and is named in
      \cref{ch:discussion} \\
    Characterise the reachable workspace & Done, and the result is negative:
      the two arms share no reachable point in front of the wearer \\
    Add two further ways of driving & Done. Both reach the arms through the
      same safety path, verified by test \\
    Identify what must be fixed before a study & Done, in priority order. The
      first item is a soldering repair \\
    \bottomrule
  \end{tabular}
\end{table}

\section{What was not done}

No participant has used the system, and no result has been validated on a
moving physical arm under load. The two-person study is designed, its
protocol is written, its analysis is coded and its ethics application is
drafted, and none of it has been submitted, because submitting a study whose
baseline condition depends on broken sensors would be dishonest.

The two-arm tasks in the study cannot currently be performed, for two separate
measured reasons: the arms share no workspace, and the object presents more
width than the gripper opens. Both are reported rather than worked around.

\section{What should happen next}

The ordering in \cref{ch:discussion} matters more than any individual item.
Repair the master. Measure where the camera is. Fix the clearance check so it
measures the whole arm rather than the centres of its joints. Put the
perception delay into the safety margin so that the margin means what it
claims. Then move the mounts and re-measure. Only then run people.

\section{Closing}

The most useful thing this project produced is probably not the system. It is
the discipline of not believing a surprising number until the instrument that
produced it has been cleared, and the record of seventeen occasions on which
that discipline was the difference between a finding and a mistake. The robot
will be rebuilt. The habit transfers.
